\scanpage{226}
dernière, on \deleted{[peut se ramener, en décomposant $Z_{s_0}$ en \deleted{\unsure{des}} \inserted{petits} chemins $a_i^0\longrightarrow a_{i+1}^0$, et transportant \deleted{\unsure{par}} horizontalement les $a_i^0$, \deleted{\unsure{et}}]} va prouver la relation plus générale
\[
(13)\qquad D_s\int_{a(s)}^{b(s)}\omega_s=\int_{a(s)}^{b(s)}(\overline D\omega)_s,
\]
valable pour des limites d'intégration $a(s),b(s)$ qui sont des \uline{sections horizontales} pour la connexion. [NB. Bien entendu, on suppose, pour donner un sens à cette formule, qu'on choisit les \unsure{$=$|mêmes} $C_s$ aux deux mb, donc (13) s'écrit en tte précision
\[
(13\ \mathrm{bis})\qquad D_s\int_{C_s}\omega_s=\int_{C_s}(\overline D\omega)_s\qquad].
\]

Pour prouver (13), on se ramène, en \inserted{se plaçant au voisinage de $s_0\in S$ et} décomposant le domaine $\int_{a(s)}^{b(s)}$, au cas où $a(s_0),b(s_0)$ sont assez voisins, de sorte qu'on peut se placer dans un ouvert $U$ où l'on a $\omega=d\varphi$. Mais alors le $1^{\mathrm{er}}$ mb est $D_s(\varphi(b(s))-\varphi(a(s)))$, le
\margin{225}
